\subsection{Custom architecture design}
\label{sec:custom-architecture}

Both custom networks take a $64\times64$ RGB image and output ten class
logits.  Each feature-extraction stage has a valid $3\times3$ convolutional
operation with stride~1 and no padding, followed by ReLU and $2\times2$ max
pooling with stride~2.  Because of this, the spatial size goes through
$64\rightarrow62\rightarrow31\rightarrow29\rightarrow14\rightarrow12
\rightarrow6$ over the three stages.  Both networks use channel widths of 32,
64 and 128 and hidden fully connected widths of 128 and 64.  The only
differences are that Model B replaces Model A's standard convolutions with
depthwise separable convolutions, and replaces its flattened head with global
average pooling (GAP).

All convolutional and dense layers have biases, and there is no batch
normalisation in either model.  During training, dropout with probability 0.2
is applied right before each hidden dense layer, and it is turned off for
validation and inference.  The final dense layer outputs raw logits, so there
is no ReLU or softmax after it.

A standard convolution has
\begin{equation}
  P_{\mathrm{conv}}=(K_hK_wC_{\mathrm{in}}+1)C_{\mathrm{out}},
\end{equation}
where the added one is the bias of each output channel.  A depthwise and
pointwise pair instead has
\begin{equation}
  P_{\mathrm{sep}}=(K_hK_w+1)C_{\mathrm{in}}
  +(C_{\mathrm{in}}+1)C_{\mathrm{out}}.
\end{equation}
A dense layer has $(N_{\mathrm{in}}+1)N_{\mathrm{out}}$ parameters.  For one
image, we count convolution MACs as
$K_hK_wC_{\mathrm{in}}C_{\mathrm{out}}H_{\mathrm{out}}W_{\mathrm{out}}$
and dense MACs as $N_{\mathrm{in}}N_{\mathrm{out}}$, and bias additions are
not counted.  Max pooling, ReLU, dropout, flatten and GAP have no trainable
parameters.  Pooling and averaging operations are also left out of the MAC
total, because the project profiler does not count them and the hand
derivation should match the profiler exactly.

\subsubsection{Model A: standard CNN}

Model A flattens the final feature map into $128\times6\times6=4{,}608$
values before the first dense layer, and
Table~\ref{tab:model-a-hand-derivation} derives the parameters and MACs of
every layer from these shapes.

\begin{table}[htbp]
  \centering
  \caption{Hand-derived parameters and MACs for Model A.}
  \label{tab:model-a-hand-derivation}
  \resizebox{\linewidth}{!}{%
  \begin{tabular}{llll}
    \hline
    Layer & Output & Parameter arithmetic & MAC arithmetic \\
    \hline
    Conv 1 & $32\times62\times62$ & $(3\times3\times3+1)\times32=896$
      & $3\times3\times3\times32\times62^2=3{,}321{,}216$ \\
    MaxPool 1 & $32\times31\times31$ & $0$ & $0$ (not profiled) \\
    Conv 2 & $64\times29\times29$ & $(3\times3\times32+1)\times64=18{,}496$
      & $3\times3\times32\times64\times29^2=15{,}501{,}312$ \\
    MaxPool 2 & $64\times14\times14$ & $0$ & $0$ (not profiled) \\
    Conv 3 & $128\times12\times12$ & $(3\times3\times64+1)\times128=73{,}856$
      & $3\times3\times64\times128\times12^2=10{,}616{,}832$ \\
    MaxPool 3 & $128\times6\times6$ & $0$ & $0$ (not profiled) \\
    Flatten & $4{,}608$ & $0$ & $0$ \\
    Dense 1 & $128$ & $(4{,}608+1)\times128=589{,}952$
      & $4{,}608\times128=589{,}824$ \\
    Dense 2 & $64$ & $(128+1)\times64=8{,}256$
      & $128\times64=8{,}192$ \\
    Output & $10$ & $(64+1)\times10=650$ & $64\times10=640$ \\
    \hline
    Total & & $692{,}106$ & $30{,}038{,}016$ \\
    \hline
  \end{tabular}%
  }
\end{table}

The convolutional stack has only $93{,}248$ parameters, while the dense head
has $598{,}858$, which is $86.5\%$ of the total.  So the head is the part that
Model B has to cut if it is to meet its parameter budget.

\subsubsection{Model B: depthwise separable CNN}

In each Model B stage, the depthwise step applies one $3\times3$ filter to
each input channel, and the pointwise $1\times1$ step then mixes those
channels.  There is only one ReLU per stage, and it comes after the pointwise
step.  Table~\ref{tab:model-b-hand-derivation} lists the two convolutions as
separate rows so that the cost of each step can be seen.

\begin{table}[htbp]
  \centering
  \caption{Hand-derived parameters and MACs for Model B. DW and PW denote
  depthwise $3\times3$ and pointwise $1\times1$ convolution.}
  \label{tab:model-b-hand-derivation}
  \resizebox{\linewidth}{!}{%
  \begin{tabular}{llll}
    \hline
    Layer & Output & Parameter arithmetic & MAC arithmetic \\
    \hline
    DW 1 & $3\times62\times62$ & $(3\times3+1)\times3=30$
      & $3\times3\times3\times62^2=103{,}788$ \\
    PW 1 & $32\times62\times62$ & $(3+1)\times32=128$
      & $3\times32\times62^2=369{,}024$ \\
    MaxPool 1 & $32\times31\times31$ & $0$ & $0$ (not profiled) \\
    DW 2 & $32\times29\times29$ & $(3\times3+1)\times32=320$
      & $3\times3\times32\times29^2=242{,}208$ \\
    PW 2 & $64\times29\times29$ & $(32+1)\times64=2{,}112$
      & $32\times64\times29^2=1{,}722{,}368$ \\
    MaxPool 2 & $64\times14\times14$ & $0$ & $0$ (not profiled) \\
    DW 3 & $64\times12\times12$ & $(3\times3+1)\times64=640$
      & $3\times3\times64\times12^2=82{,}944$ \\
    PW 3 & $128\times12\times12$ & $(64+1)\times128=8{,}320$
      & $64\times128\times12^2=1{,}179{,}648$ \\
    MaxPool 3 & $128\times6\times6$ & $0$ & $0$ (not profiled) \\
    GAP and flatten & $128$ & $0$ & $0$ (not profiled) \\
    Dense 1 & $128$ & $(128+1)\times128=16{,}512$
      & $128\times128=16{,}384$ \\
    Dense 2 & $64$ & $(128+1)\times64=8{,}256$
      & $128\times64=8{,}192$ \\
    Output & $10$ & $(64+1)\times10=650$ & $64\times10=640$ \\
    \hline
    Total & & $36{,}968$ & $3{,}725{,}196$ \\
    \hline
  \end{tabular}%
  }
\end{table}

GAP reduces the $128\times6\times6$ tensor to 128 channel means before the
dense stack, so Model B does not need the $4{,}608\rightarrow128$ layer that
dominates Model A.

The separable blocks have $11{,}550$ parameters and the GAP-based dense head
has $25{,}418$, which puts Model B $63{,}032$ parameters below the
100,000-parameter limit.  Keeping a flattened head would leave a
$589{,}952$-parameter first dense layer, and that layer alone is already over
the budget.  Narrowing the channels might also have brought the count under
the limit, but it would reduce the capacity of the feature extractor, so we
used GAP and kept the same channel widths as Model A.

\subsubsection{Compute saving and profiler cross-check}

Table~\ref{tab:separable-mac-saving} compares only the convolution
substitution, at the same input and output shapes.  Over the three stages,
depthwise separable convolution reduces the convolution MACs from
$29{,}439{,}360$ to $3{,}699{,}980$, which is a $7.96\times$ reduction.  The
reduction gets larger as the channel count grows.  For $256\rightarrow256$
channels, the exact cost per output location is
$9\times256\times256=589{,}824$ MACs for a standard convolution and
$9\times256+256\times256=67{,}840$ for the separable version, which gives an
$8.69\times$ reduction.

\begin{table}[htbp]
  \centering
  \caption{Direct MAC comparison between standard and depthwise separable
  convolution at each common stage shape.}
  \label{tab:separable-mac-saving}
  \begin{tabular}{lrrr}
    \hline
    Stage ($C_{\mathrm{in}}\rightarrow C_{\mathrm{out}}$) & Standard & Separable & Reduction \\
    \hline
    $3\rightarrow32$, $62\times62$ & $3{,}321{,}216$ & $472{,}812$ & $7.02\times$ \\
    $32\rightarrow64$, $29\times29$ & $15{,}501{,}312$ & $1{,}964{,}576$ & $7.89\times$ \\
    $64\rightarrow128$, $12\times12$ & $10{,}616{,}832$ & $1{,}262{,}592$ & $8.41\times$ \\
    \hline
    Convolution total & $29{,}439{,}360$ & $3{,}699{,}980$ & $7.96\times$ \\
    \hline
  \end{tabular}
\end{table}

When the classifier heads are included, Model B has $18.72\times$ fewer
trainable parameters and $8.06\times$ fewer MACs per image than Model A,
which are reductions of $94.66\%$ and $87.60\%$.  These are counts and not
measured speedups, because actual inference time also depends on memory
traffic, kernel support and the operations left out of the MAC count.

The forward-hook profiler gives $692{,}106$ parameters and $30{,}038{,}016$
MACs for Model A, and $36{,}968$ parameters and $3{,}725{,}196$ MACs for
Model B.  Both totals agree exactly with the hand derivations, so the tables
describe the networks that we actually trained.
